\section{Model and Opportunity}
\label{sec:model}
The measurements in \secref{sec:evidence} establish that coordination costs exist and
are substantial at small batches. They do not establish how much of that cost is
\emph{exposed}, and this distinction decides whether any redesign can help. A
synchronization stage that runs concurrently with work the consumer needs anyway costs
nothing on the critical path; the same stage serialized against that work costs its
full duration. This section builds the accounting that separates the two, then
identifies which portions in-network coordination can overlap. Ranks are allowed to
become ready at different times throughout; simultaneous readiness is a special case,
not an assumption.

\subsection{Starting from Swift's measured events}
Swift decomposes the RTT of a data packet and its acknowledgment into endpoint and fabric components~\cite{swift}. Let $t_1$ be the sender-stack transmission timestamp, $t_2$ the remote NIC receive timestamp, $t_3$ remote stack reception, $t_4$ ACK preparation, $t_5$ local NIC reception of the ACK, and $t_6$ local stack reception. The decomposition is
\begin{align}
 R_{\rm e2e} &= t_6-t_1,\\
 D_{\rm endpoint} &= (t_4-t_2)+(t_6-t_5),\\
 D_{\rm fabric} &= (t_5-t_1)-(t_4-t_2).
 \label{eq:swift}
\end{align}
The receiver returns its locally measured delay $t_4-t_2$ in the ACK. This lets the sender subtract durations without sharing a clock origin with the receiver. NIC and host clocks on each machine must nevertheless be converted to compatible units. Because transmit timestamps are stack events, unmeasured NIC transmit queuing can remain in the residual. Swift's send organization makes that component small; the formula does not turn an arbitrary application RTT into pure propagation delay.

We borrow this measurement discipline rather than Swift's congestion controller. For GPU communication, endpoint work includes GPU/NIC initiation, polling, memory visibility, and local synchronization. A GIN signal need not follow Swift's software ACK path. Similarly, a half RTT is a one-way approximation only under an explicit symmetry assumption.

\subsection{From message delay to synchronization}
For a specified message and pair of event boundaries, write
\begin{equation}
 d_{i,d}=p_i^{\rm tx}+\ell_{i,d}+p_d^{\rm rx},
 \label{eq:message}
\end{equation}
where $\ell$ contains fabric propagation, serialization, and queuing, while $p^{\rm tx}$ and $p^{\rm rx}$ contain endpoint work outside those boundaries. Serialization is charged once, according to where injection is timestamped. A collective joins these messages with a maximum over required arrivals; it does not sum all concurrent peer-to-peer traversals.

Let $r_i$ be rank $i$'s local data/resource readiness time and $\mathcal{R}_i$ its required remote peers. Both local readiness and remote announcements must hold before a baseline sender is released:
\begin{equation}
 \begin{aligned}
 h_i&=\max_{j\in\mathcal{R}_i}\{r_j+u_{j,i}\},\\
 g_i&=\max\{r_i,h_i\}+v_i.
 \end{aligned}
 \label{eq:gate}
\end{equation}
Here $u_{j,i}$ covers \ready{} initiation and propagation and $v_i$ covers the final local observation/join. For an empty peer set, take $h_i=-\infty$. In particular, early remote announcements cannot release a sender before its own $r_i$. Readiness skew is represented by $r_j$; it is not fabric queuing. In Hybrid paths, scale-up and scale-out readiness can overlap, so the network branch's latency cannot simply be added to the scale-up branch's latency.

For an EP phase, our measurement boundaries are local readiness, first \data{} injection, last \data{} injection, last required \data{} visible to its consumer, and successor release. These respectively delimit pre-synchronization, injection, tail delivery, and post-synchronization. Buffer recycling occurs after its last consumer and has a separate timestamp. They are different from Swift's six packet/ACK timestamps.

% Schematic timing of endpoint coordination and the proposed \sys{} path.
% Coordinates are schematic, not measurements. Dashed = control; solid = \data{}.
\begin{figure*}[t]
\centering
\begin{tikzpicture}[x=1cm,y=-1cm,font=\footnotesize]
\path[use as bounding box] (0,0) rectangle (17.5,9.3);
\begin{scope}[shift={(0,4.75)}]
\node[paneltitle] at (.1,.15) {(c) Aligned ranks: endpoint exchanges};
\node at (1.1,.65) {Rank 0}; \node at (5.7,.65) {Rank 1};
\draw[lifeline] (1.1,.9)--(1.1,4.2); \draw[lifeline] (5.7,.9)--(5.7,4.2);
\draw[ctrlarrow] (1.1,1.05)--(5.7,1.65);
\draw[ctrlarrow] (5.7,1.05)--(1.1,1.65);
\node[smalllabel] at (3.4,1.08) {\ready{} all-to-all};
\foreach \y in {1.8,2.04,2.28,2.95} {
\draw[dataarrow] (1.1,\y)--(5.7,{\y+.6});
\draw[dataarrow] (5.7,\y)--(1.1,{\y+.6});}
\draw[ctrlarrow] (1.1,3.62)--(5.7,4.17);
\draw[ctrlarrow] (5.7,3.62)--(1.1,4.17);
\node[smalllabel] at (3.4,3.94) {completion exchange};
\draw[<->] (6.1,1.05)--(6.1,1.65) node[midway,right,align=left] {pre\,$\simeq L$};
\draw[<->] (6.1,3.62)--(6.1,4.17) node[midway,right,align=left] {post\,$\simeq L$};
\end{scope}
\begin{scope}[shift={(8.8,4.75)}]
\node[paneltitle] at (.1,.15) {(d) Aligned ranks: \sys{} overlap};
\node at (.8,.65) {Rank 0}; \node at (3.6,.65) {\sys{}}; \node at (6.4,.65) {Rank 1};
\foreach \x in {.8,3.6,6.4} \draw[lifeline] (\x,.9)--(\x,4.2);
\draw[ctrlarrow] (.8,1.05)--(3.6,1.35); \draw[ctrlarrow] (6.4,1.05)--(3.6,1.35);
\fill[incgreen] (3.6,1.35) circle (.045);
\node[smalllabel,anchor=west] at (3.75,1.08) {all \ready{}};
\foreach \y in {1.2,1.44,1.68,2.35} {
\draw[dataarrow] (.8,\y)--(6.4,{\y+.6});
\draw[dataarrow] (6.4,\y)--(.8,{\y+.6});}
\draw[ctrlarrow] (3.6,2.72)--(6.4,3.02);
\draw[ctrlarrow] (3.6,2.72)--(.8,3.02);
\node[smalllabel,align=center] at (3.6,3.44) {\done{} follows final \data{}\\on each ordered output};
\node[anchor=west] at (.4,4.02) {No sender release wait; no new endpoint exchange.};
\end{scope}
\begin{scope}[shift={(0,0)}]
\node[paneltitle] at (.1,.15) {(a) Unequal ranks: endpoint exchanges};
\node at (1.1,.65) {Rank 0}; \node at (5.7,.65) {Rank 1};
\draw[lifeline] (1.1,.9)--(1.1,4.3); \draw[lifeline] (5.7,.9)--(5.7,4.3);
\draw[ctrlarrow] (1.1,1.0)--(5.7,1.6); \draw[ctrlarrow] (5.7,1.55)--(1.1,2.15);
\node[smalllabel] at (3.4,1.02) {Rank 1 ready later};
\foreach \y in {2.2,2.43,2.66,3.03} \draw[dataarrow] (1.1,\y)--(5.7,{\y+.6});
\foreach \y in {1.65,1.88,2.11,2.48} \draw[dataarrow] (5.7,\y)--(1.1,{\y+.6});
\draw[ctrlarrow] (1.1,3.2)--(5.7,3.8);
\draw[ctrlarrow] (5.7,3.72)--(1.1,4.32);
\node[smalllabel,align=center] at (3.4,4.15) {Readiness and completion\\still traverse between endpoints};
\end{scope}
\begin{scope}[shift={(8.8,0)}]
\node[paneltitle] at (.1,.15) {(b) Unequal ranks: bounded staging};
\node at (.8,.65) {Rank 0}; \node at (3.6,.65) {\sys{}}; \node at (6.4,.65) {Rank 1};
\foreach \x in {.8,3.6,6.4} \draw[lifeline] (\x,.9)--(\x,4.3);
\fill[pattern=north east lines,pattern color=black!35] (3.43,1.43) rectangle (3.77,1.87);
\draw[ctrlarrow] (.8,1.0)--(3.6,1.3);
\draw[ctrlarrow] (6.4,1.55)--(3.6,1.85);
\draw[dataarrow] (.8,1.13)--(3.6,1.43);
\draw[dataarrow] (.8,1.36)--(3.6,1.66);
\draw[dataarrow] (.8,1.59)--(3.6,1.89);
\draw[dataarrow] (3.6,1.9)--(6.4,2.2);
\draw[dataarrow] (3.6,2.13)--(6.4,2.43);
\draw[dataarrow] (3.6,2.36)--(6.4,2.66);
\draw[dataarrow] (.8,2.18)--(3.6,2.48)--(6.4,2.95);
\foreach \y in {1.68,1.91,2.14,2.51} \draw[dataarrow] (6.4,\y)--(.8,{\y+.6});
\draw[ctrlarrow] (3.6,2.72)--(6.4,3.02);
\draw[ctrlarrow] (3.6,2.88)--(.8,3.18);
\node[smalllabel,anchor=east] at (3.2,1.68) {buffer};
\node[smalllabel,align=center] at (3.6,3.6) {Wait for \ready{} at the \sys{} node;\\preserve output service and ordering};
\node[anchor=west] at (.4,4.18) {Skew coverage is not an additional fixed saving.};
\end{scope}
\end{tikzpicture}
\caption{Idealized projection of the two-boundary design, not a measured Direct trace. Solid arrows carry \data{}; dashed arrows carry control. The left-hand reference assumes an exposed tail-signal interval of approximately $L$ after data arrival; actual Direct timing follows Equation~\eqref{eq:tailgain}. \sys{} overlaps readiness with upload and appends ordered \done{} at each destination's tail. Early data is buffered under skew. Required local completion and buffer ownership are retained.}
\Description{Four timing panels compare endpoint and network coordination. The top pair shows unequal readiness and buffering; the lower pair shows aligned readiness as an ideal reference. Solid arrows are data and dashed arrows are control.}
\label{fig:timing}
\end{figure*}


\subsection{Direct's tail: initiation versus visibility}
\label{sec:directtail}
In cross-node Direct, the tail barrier waits for applicable asynchronous memory operations, flushes the GIN queues, and joins local workers before issuing the group-wide completion signals~\cite{deepepcode}. The GIN flush contract guarantees that PUT source buffers are safe to reuse, but does not identify when the destination GPU observes the data~\cite{ginheader,gindocs}. The GDAKI implementation performs a send-completion-queue wait~\cite{gdakicode}. Sender completion and remote visibility are therefore represented by distinct timestamps below.

Let $c_j$ be the time sender $j$ completes the transfer-progress waits and local join that precede tail signaling. Let $w_{j,d}$ cover the subsequent signal initiation and propagation to $d$, and $v_d$ its final observation overhead. Tail control releases $d$ at
\begin{equation}
 z_d=\max_{j\in\mathcal{R}_d\cup\{d\}}(c_j+w_{j,d})+v_d.
 \label{eq:tailsignal}
\end{equation}
The self term represents local completion, with no propagation over the network. Let $\tau_d$ be the time all required input data becomes visible at $d$. At a comparison point where other local dependencies are unchanged, the baseline cannot consume before $\max(\tau_d,z_d)$. \sys{} can instead release on its ordered output notification at $\tau'_d+\eta_d$, where $\eta_d$ is the notification lag after data visibility. Thus,
\begin{equation}
 \Delta T_{\rm tail,d}=\max(\tau_d,z_d)-(\tau'_d+\eta_d).
 \label{eq:tailgain}
\end{equation}
For a tail-only comparison with identical payload arrival we have $\tau'_d=\tau_d$, so
the saving is the exposed control wait, $\max(0,z_d-\tau_d)$, less the lag. This
describes the Direct optimization target without assuming that a sender's flush return
coincides with a receiver's arrival time. The sender retains any source-lifetime
completion obligation; the design decouples that obligation from target notification
through generation ownership.

Equation~\eqref{eq:tailgain} also explains what our measurements can and cannot settle.
The signal window of \secref{sec:paired} bounds $z_d - c_j$, the interval from signal
initiation to observation, and shows it is 11.6--14.3\us{} at EP8. It does not reveal
$z_d - \tau_d$, the exposed wait, because $\tau_d$ was not observed: doing so requires
timestamping destination visibility separately from sender completion on the same
configuration. The stage's existence is established by the source; its exposed fraction
remains open, and we do not assume it equals $L$.

\subsection{Unequal arrivals and finite service}
\label{sec:skew}
With admitted staging capacity, a \sys{} sender can start uploading at $s_i^{\rm up}=r_i+\delta_i$, where $\delta_i$ is its announcement and injection-preparation delay. If $a_i$ includes \ready{} initiation and propagation from local readiness to the \sys{} node, the network gate opens at $o=\max_i(r_i+a_i)$. Local readiness thus controls upload, while collective readiness controls network delivery. For destination $d$, let $\lambda_d$ be the time its epoch and any required layout become usable. Its release condition is $o_d=\max(o,\lambda_d)$.

A small service model makes the role of buffering explicit. Choose an ordered sequence of non-overlapping data units for one output. Let $A_k$ be unit $k$'s arrival time at the \sys{} node, $s_k$ its output service time, and $E_k$ its output completion time. Under a non-preemptive output model,
\begin{equation}
 E_k=\max\{A_k,o_d,E_{k-1}\}+s_k.
 \label{eq:queue}
\end{equation}
Downstream propagation and visibility determine when this output becomes consumable. The units can be packets or explicitly modeled fragments; a cut-through implementation needs compatible boundaries. A finite staging pool also constrains arrivals through backpressure. Applying the same service accounting to the baseline, with injections delayed until $g_i$, produces a comparable completion time.

Early data can accumulate while a late rank is unready, but the gate does not erase that late rank or drain the queue instantaneously. Early-upload coverage and whole-operation latency reduction are different quantities. In particular, an egress bottleneck can consume the head start. Readiness skew therefore cannot be added as an independent bonus to Equation~\eqref{eq:gain}. This service model also explains why a count of admitted bytes alone is insufficient evidence of latency savings.

\subsection{Two one-way traversals in the ideal case}
\label{sec:ideal}
Figure~\ref{fig:timing} projects the two coordination boundaries onto a simple aligned schedule. Let $L$ denote one-way propagation after removing endpoint overheads, $S$ the common payload injection duration, and $P$ common retained processing. The symmetric fabric RTT is $R_f=2L$. For the tail, additionally assume the exposed signal interval in Equation~\eqref{eq:tailgain} is approximately $L$. This is a reference operating point, not a consequence of the GIN flush contract. With negligible control-packet serialization,
\begin{equation}
 T_{\rm base}=P+L+(S+L)+L=P+S+3L.
 \label{eq:base}
\end{equation}
Suppose sender-to-node propagation is $a$, node-to-receiver propagation is $b$, and $a+b=L$. In the symmetric reference all \ready{} messages reach the \sys{} node after $a$. \data{} is already following them. Once ready, the node streams at the reference service rate, giving payload completion at $S+L$ rather than $L+S+L$. The midpoint location is illustrative: equal upstream readiness distances and comparable service, not geometric midpoint placement alone, establish this schedule.

At the tail, the node appends \done{} after the destination's final data on an ordered delivery channel. Data and notification share the remaining downlink. The consumer does not initiate another inter-endpoint exchange after receiving its payload. Let $\epsilon$ include the incremental cost of announcements, staging, ordering, notification, and changed queuing. Then
\begin{align}
 T_{\text{\sys}}&=P+S+L+\epsilon,\\
 \Delta T&=2L-\epsilon=R_f-\epsilon.
 \label{eq:gain}
\end{align}
Each eligible boundary contributes approximately $L$ while preserving the complete data path and local input dependency. A successor requiring global consumer completion needs an additional join outside this model.

For a path with $k\in\{0,1,2\}$ eligible boundaries, the same reference gives
\begin{equation}
 \Delta T_k\simeq kL-\epsilon_k.
 \label{eq:eligible}
\end{equation}
This is neither a measured speedup nor a bound for arbitrary queueing schedules, and
$L$ is not directly observable from our data: the 15\us{} entry window of
\secref{sec:paired} contains GPU-side initiation, polling, and peer-readiness skew, so
substituting it for $L$ would overstate the fabric term substantially.
Section~\ref{sec:sensitivity} therefore treats $\epsilon_k$ as a swept parameter rather
than a fitted one.
